\input{../tex-inputs/latex-leseansicht-vorspann}

\section[Arthur Schnitzler an Gustav Schwarzkopf, 17. 4. 1928]{L04461 Arthur Schnitzler an Gustav Schwarzkopf, 17. 4. 1928}
\nopagebreak\mylabel{L04461v}
\rehead{ }\normalsize\beginnumbering\briefempfaengerindex{Schwarzkopf, Gustav@\textsc{Schwarzkopf, Gustav}!zzzSchnitzler, Arthur@\emph{von Arthur Schnitzler}!1928-04-171@{17. 4. 1928}|(be}
\toendnotes[C]{\smallbreak\pagebreak[2]}
\correspDesc{Versand  durch Arthur Schnitzler am 17. 4. 1928 in Mittelmeer
\newline{}Übermittlung  am 18. 4. 1928 in Piraeus
\newline{}Erhalt  durch Gustav Schwarzkopf im Zeitraum [18. 4. 1928 – 22. 4. 1928?] in Wien}\toendnotes[C]{\smallbreak}
\Standort{DLA, A:Schnitzler, HS.1985.1.1897.}
\physDesc{Bildpostkarte, 182 Zeichen
\newline{}Handschrift: Bleistift, deutsche Kurrent
\newline{}Versand: Stempel: »\nobreak{}\oindex{Piraeus@\textbf{Piraeus}|pwk}\griechisch{Πειραιευς}, 18. \griechisch{απρ.} 28., 3 \griechisch{ε}, \griechisch{αποστολη}\nobreak{}«.  }\toendnotes[C]{\smallbreak}\pstart{}{\pb}Austria\oindex{Österreich@\textbf{Österreich}|pw}\pend{}\pstart{}Hr Gustav Schwarzkopf\pend{}\pstart{}Wien I\oindex{I., Innere Stadt@\textbf{I., Innere Stadt}|pw}\pend{}\pstart{}Tiefer Graben 17\oindex{Wien@\textbf{Wien}!I., Innere Stadt@\textbf{I., Innere Stadt}!Tiefer Graben 17@\textbf{Tiefer Graben 17}|pw}\pend{}{\bigskip}
\pstart
           \noindent{}{\pb}\textcolor{gray}{\textbf{COSULICH LINE\orgindex{Austro-Americana@Austro-Americana|pw}}}{\\}\textcolor{gray}{\textbf{S. S. »STELLA D’ ITALIA«}}\orgindex{Stella d’Italia@Stella d’Italia|pw}\pend
           \vspace{1em}
\pstart
           \raggedleft{}{\pb}17/4 28\pend
           \vspace{0.5em}
\pstart
           Abends, nach Abfahrt
               von Katakolon\oindex{Katákolo@\textbf{Katákolo}|pw}, viele 
               herzliche Grüße,
      auch von meiner lieben
               Reisebegleiterin\pwindex{Cappellini, Lili 13.\,9.\,1909 Wien – 26.\,7.\,1928 Venedig@\textsc{Cappellini, Lili} (13.\,9.\,1909 Wien – 26.\,7.\,1928 Venedig)|pwv}. Grüßen Sie
      Dr Max\pwindex{Schwarzkopf, Max 12.\,6.\,1857 Wien – 14.\,4.\,1928 ebd.@\textsc{Schwarzkopf, Max} (12.\,6.\,1857 Wien – 14.\,4.\,1928 ebd.)|pw}.\pend
           \pstart Ihr \spacefill\mbox{Arthur}\pend{}\selectlanguage{ngerman}\endnumbering\briefempfaengerindex{Schwarzkopf, Gustav@\textsc{Schwarzkopf, Gustav}!zzzSchnitzler, Arthur@\emph{von Arthur Schnitzler}!1928-04-171@{17. 4. 1928}|)be}\mylabel{L04461h}
\begin{anhang}
\end{anhang}\newcommand{\dateiname}{L04461}\newcommand{\titel}{Arthur Schnitzler an Gustav Schwarzkopf, 17. 4. 1928}\newcommand{\editorInnen}{Herausgegeben von Jahnke, SelmaMüller, Martin Anton}\input{../tex-inputs/latex-leseansicht-abspann}
